\documentclass[11pt,a4paper]{article}
\usepackage[margin=23mm]{geometry}
\usepackage[T1]{fontenc}
\usepackage{lmodern,amsmath,amssymb,booktabs,longtable,array,enumitem}
\usepackage[hidelinks,bookmarks=false]{hyperref}
\emergencystretch=3em
\newcommand{\C}{\textbf{(c)}\ }
\newcommand{\D}{\textbf{(d)}\ }
\newcommand{\file}[1]{\nolinkurl{#1}}
\newcommand{\test}[1]{\textit{Independent check:} \file{review_draft3_checks.py}, group #1.}
\newcolumntype{L}[1]{>{\raggedright\arraybackslash}p{#1}}
\title{Independent cross-review of paper draft 3}
\author{Review for Gustav Delius}
\date{4 November 2026}
\begin{document}
\maketitle

\begin{abstract}
\C The draft incorporates most audited corrections: the bad-equation
one-body matching, corrected Wiener--Hopf normalisation, finite-sea endpoint
identities, Gamma candidate poles and conditional ring algebra survive.
There are \textbf{11 required corrections} below. The principal failures are
the four-term Krein source and its claimed physical consequences, the signs
of the collective symplectic matrices and definitions of projector sources,
and inconsistent one-loop units. Passing paper scripts do not certify those
statements. \D Acceptance as an audited consolidation should await the
corrections; neither an exact S-matrix nor a physical projected determinant
is established by this review.
\end{abstract}

\section{Scope, conventions and evidence}
\C Reviewed source: \file{paper/boussinesq_soliton_scattering.tex}, draft 3,
including the coordinator's revised Section 10. The source SHA-256 is
\begin{quote}\small\ttfamily
81d98ad5a437ab876df7f6b359ffd9c5800\newline
df84bd5c9087238bd19b8432f2d03
\end{quote}
All line locators below refer to that snapshot. The required
\file{git log --format='--- \%h \%s\%n\%b' a111f8c^..HEAD}
digests were read, including the later warning in commit \file{38f5b2d}.
No pull, index update, commit, stash or branch change was performed.

Labels mean (a) literature proof, (b) literature heuristic, (c) project
derivation or reproducible source observation, and (d) unresolved physical
identification or research judgment. All source-comparison findings and table
entries below are (c), with physical limitations explicitly (d).
No external paper or abstract was newly read, no external bibliography was
reverified, and no new arXiv ID or DOI is asserted. Literature names refer
to the formulas and attribution already recorded in the local packages.
The controlling internal texts are the relevant sections of
\file{phase2/audit_defq.tex}, \file{audit_det_projected.tex},
\file{audit_projected_measure.tex}, \file{audit_toda_dictionary.tex},
\file{audit_frame_dictionary.tex}, \file{audit_phase2i.tex},
\file{audit_onebody.tex}, \file{audit_onebody_scheme.tex},
\file{audit_smx.tex}, their packages, \file{phase2/errata.tex},
\file{phase2/twobody_toda.tex}, and \file{phase2/report.tex}.
Audit verdicts override package prose; the later explicit reference-term
calculation also requires reconsidering the older frame audit.

\C The independent checker uses exact SymPy and 60-digit mpmath, imports
no project implementation, and writes no files. Its group names G, E, W,
K, F, C, S, R, D and I are used below. Source checks are distinguished from
fresh computations. The four pre-existing \file{paper/scripts/*.py} files
were all executed; the new review checker is also executed. Phase-2 scripts
were inspected, not indiscriminately rerun, since many write existing data.
Only the requested three deliverables are added to the repository.

\section{Required corrections}
\begin{enumerate}[label=\textbf{R\arabic*.},leftmargin=*,itemsep=1.2em]

\item \textbf{Replace the four-term Krein source; weaken the physical exclusion.}
\C \textit{Location:} Section 8.3, lines 805--821; summary item 8;
Appendix A's antisymmetric-function row; the rank-two route in Section 10(2).
The equation imported from \file{audit_frame_dictionary.tex} is wrong.
Use notation that avoids the draft's overloaded $f$: let
$B=F'$, $b=B'$, $c=\sqrt\pi$, and, for $u>v>0$,
\begin{align}
 H(u,v)&=2\int_0^v b(u-w)B(v-w)\,dw,\qquad
 \widehat H(u,v)=H(u,v)-2c[B(v)-c],\label{rev:H}\\
 (\partial_u+\partial_v)H&=2b(u)B(v),\\
 (\partial_u+\partial_v)\widehat H&=2b(u)B(v)-2c\,b(v).
 \label{rev:source}
\end{align}
The old source, even granting its intended derivative convention,
$2b(u)[B(v)-c]$, differs by $2c[b(u)-b(v)]$.
At $(u,v)=(4,1)$ the corrected source is
$-0.3471042546629204553\ldots$, versus
$+0.8377480606920601269\ldots$ for the old source.
\textit{Evidence:} \file{twobody_toda.tex}, equation \file{eq:PDE};
fresh transform inversion and differentiated-convolution bookkeeping.
\test{K}

\C This repair does \emph{not} produce an arbitrary freedom in the
specified microscopic Toda phase. With its source and boundary data fixed,
the difference of two solutions obeys
$(\partial_u+\partial_v)\Delta=0$ and vanishes at $v=0$;
characteristics give $\Delta=0$ in $u>v>0$.
The correct test of $A_\lambda(y_1,y_2)$ is
\[
 \left[\frac{B(u_1)}{2\sqrt{y_1}}\partial_{y_1}
       +\frac{B(u_2)}{2\sqrt{y_2}}\partial_{y_2}\right]A_\lambda,
\]
not $(\partial_{y_1}+\partial_{y_2})A_\lambda$.
The independent checker finds this weighted derivative nonzero at the
actual WH coordinates corresponding to $(u_1,u_2)=(4,1)$.

\C The nonzero $3\times3$ minor only excludes Schmidt rank at most two;
it does not by itself prove infinite rank or restrict every admissible
physical phase to a bilinear. Infinite rank of this particular rational
function can instead be proved by its distinct moving poles at $y_1=-y_2$.
The evenness in $r$ and poles at $r=\pm i$ are algebraic facts.
\D Universal physical parity, a pole-free complex domain for the unknown
quantum frame, and restriction of every frame correction to the specified
one-body bilinear have not been established by those algebraic facts.
The later directional residual even contains an additive $\pi$; its physical
normalisation is explicitly open.

\textit{Fix:} retain (c) for the counterexample to reciprocity alone and for
uniqueness of the \emph{specified} Toda solution conditional on the corrected
edge equation. Restrict the Schmidt statement to the stated finite-rank
frame ansatz. Withdraw the blanket three-constraint physical exclusion and
the inference that only a rank-two continuum operator problem remains.
Keep the general continuum quantum map and physical phase (d).

\item \textbf{Correct both symplectic matrices and the projector inputs.}
\C \textit{Location:} Section 7.1, \file{eq:W3W4}, lines 605--613.
For the printed order $(Z_v,Z_X,Z_C,Z_r)$ and
$\Omega=\int d\pi\wedge d\phi$, the audited matrices are
\[
 W_3=\begin{pmatrix}0&P_v&\Pi_v\\-P_v&0&0\\-\Pi_v&0&0\end{pmatrix},\qquad
 W_4=\begin{pmatrix}0&P_v&\Pi_v&0\\-P_v&0&0&0\\-\Pi_v&0&0&-1\\0&0&1&0\end{pmatrix}.
\]
Both printed matrices are their negatives. Rank, null vector and
$\det W_4=P_v^2$ are unchanged, which explains why paper check 2b passes.
The projector tail expression is correct with the \emph{corrected} matrix
and
\[
 b_A=\Omega(Z_A,\eta),\qquad
 b_r=\int w\eta_\phi\,d\xi,\qquad
 b_X=\int(vu_0'\eta_\phi+u_0\eta_\pi)\,d\xi.
\]
The paper instead identifies $b_r$ with $\int\eta_\pi$ and $b_X$ with
$\Omega(Z_v,\eta)$, both incorrectly. For $\eta=(w,0)$ its proposed $b_r$
vanishes, whereas the actual constant tail is $\int w^2d\xi=3k/7$.
\textit{Evidence:} \file{audit_det_projected.tex}, equations
\file{eq:W3_audit}, \file{eq:Winv_tails}; direct field-pairing quadrature
at both velocity signs. \textit{Fix:} use the displayed corrected matrices
and component definitions; retain the endpoint obstruction. \test{G}

\item \textbf{Use one loop-order convention and fix the finite-tadpole sign.}
\C \textit{Location:} Section 7.1, lines 622--641, especially
\file{eq:E1Pinv}, \file{eq:Fg} and assumptions (A),(B).
The paragraph defines $E=E_0+\hbar E_1|_P+\cdots$, so $E_1|_P$ is a
coefficient. It then divides that coefficient and $L_1$ by $\hbar_B$ in
\file{eq:Fg}. The cutoff expression similarly adds $3C_\Lambda M$, which
contains $\hbar$, to a trace and cutoff endpoint written without $\hbar$.
\textit{Fix:} consistently write
\[
 E(P)=E_0(P)+\hbar_B F(P)+\cdots,\quad
 L=L_0+\hbar_B\ell_1+\cdots,\quad F=e-vp=-\ell_1,
\]
with $e,p$ coefficient shifts and counterterm $3C_\Lambda M/\hbar_B$ in
the coefficient trace. Alternatively put $\hbar_B$ on every physical
trace and endpoint and call $\Delta E_1=\hbar_BF$ the correction.
The analytic right-hand side of \file{eq:Fg} itself is correct.

\C Adding $+2\hbar_Bc_{\rm fin}\sin\alpha$ to the physical $L_1$ changes
$F$ by $-2c_{\rm fin}\sin\alpha$, so the leading change in $F/|P_0|$ is
$-3c_{\rm fin}/\alpha^2$, not the printed positive sign.
The audit's positive sign is for a term added to the \emph{energy}.
\textit{Evidence:} \file{audit_det_projected.tex}, its finite-subtraction
paragraph, and \file{audit_onebody_scheme.tex}'s explicit units. \test{E}

\item \textbf{Differentiate the contact term in the paper's notation.}
\C \textit{Location:} Section 8.2, line 795, immediately before
\file{eq:R0c1}. The paper defines $f=F'=B$ in \file{eq:WHlap} and
$\psi_-=f'/2$, but copies the audit's contact term while the audit uses
$f=G''=B'$. With the paper's notation the contact term in
$\partial_{u_2}\Phi_q$ is $2\sqrt\pi f'(u_1-u_2)$, equivalently
$4\sqrt\pi\psi_-(u_1-u_2)$, not $2\sqrt\pi f(u_1-u_2)$.
\textit{Evidence:} Leibniz differentiation of \eqref{rev:H} and
\file{audit_phase2i.tex}, ``double Laplace formula'' paragraph.
\textit{Fix:} change that derivative; retain the corrected double-transform
prefactor $2\pi$ and its contact prescription. \test{K}

\item \textbf{Repair the tau-zero witness, retaining the no-breather conclusion.}
\C \textit{Location:} Section 9.2, line 856 after \file{eq:Agoodbr}.
At $\cos I=-1$, $X=1/\sqrt A$ gives
$\tau=2-2/\sqrt A$, not zero for $0<A<1$.
Also the assertion of a positive root for \emph{every} $\cos I$ when
$A\le0$ fails at $A=0$, $\cos I\ge0$.
\textit{Fix:} choose $\cos I=-1$ and
\[
 X=\frac1{1+\sqrt{1-A}}\quad(A<1).
\]
This positive root covers negative $A$, $A=0$ and $0<A<1$ uniformly.
The inequality $A<1/4$ and singularity of the specified good Hirota family
survive; no general quantum-bound-state exclusion follows.
\textit{Evidence:} direct polynomial substitution; controlling
\file{audit_smx.tex} regularity argument. \test{S}

\item \textbf{Make quantum crossing and the Gamma ansatz restrictions explicit.}
\C \textit{Location:} Section 9.1, lines 828--834, and Section 9.4.
The classical crossing-type identity in Proposition 4.2 does not establish
that an exact nonrelativistic quantum amplitude ``must'' obey the same
relation. \file{audit_smx.tex} stipulates crossing on compatible branches;
it does not prove it for the unknown physical S-matrix.
Likewise Stirling matching does not force \emph{any} single Gamma quotient
to have the paper's argument if arbitrary two-body exponential prefactors
are allowed. A compensating leading phase can cancel the Stirling phase
for another real exchange-odd argument. The root-difference rejection
script tests a particular finite basis, not that unrestricted theorem.
\textit{Fix:} label the exact crossing postulate (d), and the ensuing
family algebra (c) conditional on it. State the admissible prefactor class
before a necessity claim; otherwise present the spectral Gamma quotient
as the tested candidate. Preserve the correctly restricted mirror and
nine-ratio dilogarithm exclusions. \test{S; the restricted residue proof
is also checked in the existing paper script and \file{audit_smx.py}}

\item \textbf{Restore equation attribution and physical qualifications in summaries.}
\C \textit{Location:} summary item 7 (line 88), abstract, Section 5.7,
and Section 10's introduction and items 3--4.
The summary's \emph{both} tag includes $M(v)=2\sqrt{1-v^2}$ and the
good-equation geometry; the bad profile instead has
$M(v)=2\sqrt{v^2-1}$. The value $C-s_1=-1/(4\pi)$ is for Toda/continued
bad, whereas good gives $+1/(4\pi)$. The detailed Section 7 distinction
is correct and must govern the summary. The abstract and summary item 5
must also say that the \emph{bulk-plus-defect totals} are conserved;
localised defect charges alone change during absorption, as their nonzero
vacuum jumps in Section 5.7 already show.

\C Zero counted real growing modes does not prove complete linear
stability: the controlling defect audit leaves off-axis completeness open.
The abstract's ``strictly linearly stable'', Section 5.7's ``unique stable
vacuum'', and Section 10(4)'s unqualified stability phrase exceed that
status. The no-stable--stable-edge obstruction survives because every
other vertex has a demonstrated growing mode.
\D Keep the stronger spectral statement conditional on completeness.
The possible \emph{quantum} inelastic channels in Section 10(3) are a
project open question (d), not a literature heuristic (b) established by
the locally recorded classical CZ result.

\C Finally, Section 10's claim that the counterterm was traced to a
specific short-wave band repeats the report's motivation more strongly
than the vacuum-matching derivation supports. The scheme audit derives a
homogeneous elastic difference; \file{good_renorm.tex} explicitly warns
against inferring whole-trace support from light-root cancellation.
\textit{Fix:} retain the PI's ranking but describe that band attribution
as motivation, not an additional demonstrated mechanism. State the
one-body and WH results as conditional traces/edge limits, and the good
acoustic scheme as an available thermodynamic condition, not the selected
physical prescription. \test{E; G; D. Completeness is not asserted.}

\item \textbf{Specify the time-dependent two-soliton fluctuation operator.}
\C \textit{Location:} Section 10(2), line 902, referring to
``the fourth-order operator $\mathcal D_v$'' as the two-soliton operator.
Section 7.1 defines $\mathcal D_v$ around a single stationary travelling
profile in its co-moving frame. Two distinct velocities have no common
stationary frame. Around the good two-soliton field $u_2(x,t)$ the lab-frame
operator is
\[
 \mathcal D_{[u_2]}=\partial_t^2+\partial_x^4-\partial_x^2
                      +6\partial_x u_2(x,t)\partial_x.
\]
\textit{Fix:} use this operator (or its fully time-dependent transform),
with one-body $\mathcal D_v$ only as asymptotic input. This is not a choice
of determinant prescription. \test{G, independent quadratic-action variation}

\item \textbf{Define the auxiliary functions and the residual being expanded.}
\C \textit{Location:} \file{eq:defcharges}, Section 5.7, and
\file{eq:R0c1}, Section 8.2. The body never defines the defect polynomials
$\rho(\sigma),\tau(\sigma),h(\sigma)$ and their integration constants;
the isolated development note in Appendix A is insufficient.
Using $A=\eta/\sqrt3$ (unrelated to the interaction constant), supply
\[
 \rho(\sigma)=\frac{\sigma-\sigma^3/3-s+s^3/3}{2A},\qquad
 \tau'(\sigma)=-\sigma\rho(\sigma),\quad
 h'(\sigma)=-\rho(\sigma),\quad \tau(s)=h(s)=0.
\]
These are the audited reference and gauge conventions, not new choices.
\textit{Evidence:} \file{audit_defq.tex}, ``Classical charges rederived''.
\test{D}

\C Define the edge residual in the paper before using $R_0(Y,r)$:
subtract $H_{\rm cl}=4rY+2Y(1-r^2)\log[(1-r)/(1+r)]$ from
$H(u(Y),u(r^2Y))$, and identify $c_1$ as its formal $1/Y$ coefficient.
This makes explicit that it is the \emph{two-term directional} residual,
not the four-term phase or a physically normalised scalar phase.
\textit{Evidence:} \file{audit_phase2i.tex}, saved-data definitions;
\file{twobody_toda.tex}. \test{C, formal coefficient integrals;
the asymptotic physical identification remains (d)}

\item \textbf{Replace Appendix A's blanket formula-coverage claim with exact scopes.}
\C \textit{Location:} Appendix A, line 929 and the new quantum rows.
All 64 distinct cited script paths exist, but the opening claim that the
paper scripts assert every new formula is false. For example:
\begin{itemize}[nosep]
\item draft-3 check 2b verifies determinants/moments, not field pairings,
and accepts the wrong sign matrices; it does not test the stated projector;
\item 4a's reciprocity test is an algebraic rewrite of the endpoint product,
not a comparison with independently computed forward/reverse phases;
\item 4d only tests minors of an assumed expression; it does not construct
the four-term phase or test its actual PDE or physical parity;
\item 5e asserts cancellation of scalar symbols after assuming a crossing
law; it does not integrate $J$; 5d inserts the desired derivative identity;
\item 1b does not construct the sewing matrix; 1c does not independently
derive the quadratic energy operator; 5a does not compute a Gamma residue.
\end{itemize}
\C These are limits of the assigned checks, not reasons to discard
independent audited results. Name the controlling audit and actual test
group beside each claim. Direct field pairings, actual finite-sea reciprocity,
Gamma residues and independently differentiated ring phases are added in
this review. \test{G, F, S, R, I; source-inspection ledger below}

\item \textbf{Update stale failure and script-status statements.}
\C \textit{Location:} Appendix A's opening, Hirota row and failures item 2;
the \file{audit_onebody.py} citations in Section 7.2 and Appendix A.
The current Hirota script no longer contains an assertion bypassed by
\file{or True}; the repair is recorded in \file{errata.tex}, N05, and
\file{scripts_inventory.md}. Record that repair as historical, not current.
\file{audit_onebody.py} still fails its unchanged line-383 tolerance:
the saved discrepancy is $2.2546734882\times10^{-4}>3\times10^{-6}$.
The later audit diagnoses truncation error in a true limiting identity;
it does not make the old checker pass. Explicitly annotate its failed-run
status and use \file{audit_onebody_scheme.py} for the surviving conclusion.
Also classify \file{phase1/scripts/figures.py} as a generator, not a
verification script. \file{toda_scaling_fast.py} gates its printed marker
with a Boolean but does not raise on failure; a zero exit status alone is
not evidence. \textit{Fix:} distinguish available source, asserted coverage,
known historical failure, and actually rerun success. \test{I}
\end{enumerate}

\section{Display-by-display cross-check}
\C The following ledger covers all \textbf{38 labelled displayed equations}
in Sections 5.7, 7, 8 and 9. Labels are source labels, not newly assigned
equation numbers. ``Agrees'' means comparison with the specified controlling
audit or the fresh check, under the qualifications stated here. It is not
an independent derivation of the full quantum theory. Internal audit names
below abbreviate \file{phase2/audit_*.tex}; sources of new computations are
the indicated groups in the review checker and the executed paper checkers.
\small
\begin{longtable}{@{}L{3.0cm}L{1.6cm}L{10.6cm}@{}}
\toprule Paper equation label & Equation & Verdict, evidence and scope\\\midrule\endhead
\file{eq:defcharges} & good & Agrees with defq, action-derived charges; add primitives (R9). Paper check 1a verifies vacuum jumps, not all conservation laws.\\
\file{eq:detMdef} & good & Agrees with defq/frames\_defects, including prefactor and effective spectral parameter. D rechecks thresholds.\\
\file{eq:PD2trace} & good & Agrees with defq: correct momentum numerator and positive symplectic-energy sign; continuum contractions, not complete quantum vacua.\\
\file{eq:defUVct} & good & Agrees with defq. D rechecks decomposition; finite boundary subtraction remains open.\\
\file{eq:Kvfact} & good & Agrees with det\_projected; executed paper 2a checks factorisation and modes.\\
\file{eq:W3W4} & good & Wrong overall signs, though determinant/rank survive (R2); G checks field pairings.\\
\file{eq:boxnorm} & good & Agrees with det\_projected, including finite constant; executed paper 2c checks quadrature.\\
\file{eq:E1Pinv} & good & Correct coefficient-level Legendre identity; inconsistent later units (R3).\\
\file{eq:Fg} & good & Analytic trace and series agree; left-hand normalisation needs R3. Conditional determinant identification. E.\\
\file{eq:CEhdef} & Toda/bad & Agrees with onebody\_scheme and precision package's classical dictionary; $P_0$ is Toda momentum here.\\
\file{eq:CTnum} & Toda/bad & Central value agrees. Quoted envelope is sensitivity, not a rigorous error bound; see optional rounding improvement.\\
\file{eq:CvCE} & Toda/bad & Agrees with onebody\_scheme, conditional WH limit; F rechecks finite-sea identities, E the sound shift.\\
\file{eq:CTdecomp} & Toda/bad & Correct $1/(2\pi)$ conditional edge coefficient; not a fixed-nonzero-momentum limit theorem. E/W.\\
\file{eq:Fac} & bad & Algebraic continuation and $2/\pi$ correct. Physical unstable-mode determinant remains (d). E.\\
\file{eq:Wg} & good & Agrees with onebody\_scheme; thermodynamic vacuum coefficient. E.\\
\file{eq:Wbac} & bad & Correct formally continued strain energy and sound shift; physical bad vacuum qualified. E.\\
\file{eq:WTB} & Toda/bad & Correct harmonic-lattice vacuum energy in stated classical units. E.\\
\file{eq:DeltaHmatch} & bad & Correct homogeneous elastic matching, no soliton fit. Uniqueness is for the quadratic local coefficient. E.\\
\file{eq:Cmatched} & bad & Correct conditional trace coefficient $1/(2\pi)$. E.\\
\file{eq:schemeinv} & bad & Correct $-1/(4\pi)$ under common finite quadratic subtractions. E.\\
\file{eq:Cgonshell} & good & Correct conditional $+1/(4\pi)$ in the stated acoustic scheme, not a forced scheme. E.\\
\file{eq:endids} & Toda & Agrees with toda\_dictionary/frame\_dictionary; F solves fresh Fredholm equations.\\
\file{eq:yjac} & Toda & Correct factors of $a,h$, including energy Jacobian; fresh inverted-momentum differentiation in F.\\
\file{eq:MQresp} & Toda & Correct endpoint response and thermodynamic denominator, with edge subtraction; audited pressure response. F.\\
\file{eq:recip4} & Toda & Correct four-term identity; F compares actual directional phases, not an assumed product.\\
\file{eq:S0D0} & Toda & Agrees with frame\_dictionary, retaining endpoint terms and $a=1$. Paper 4b tests rectangular difference and large-$Y$ limit.\\
\file{eq:WHsym} & Toda & Correct nonzero-frequency symbol and factor product. W. Zero-frequency matching not proved.\\
\file{eq:WHlap} & Toda & Correct $\sqrt\pi$ momentum/derivative transforms; half-line matching conditional. W.\\
\file{eq:duasymp} & Toda & Correct $c_{\rm edge},d_0,e_0$ under transform input; phase2i symbolic reversion and paper 4c, independently inverted in W.\\
\file{eq:R0c1} & Toda & Formal integrals agree with phase2i and fresh C. Define residual (R9); uniform integrated asymptotics remain conditional.\\
\file{eq:Alambda} & Toda & Correct rational counterexample and scaling; claimed physical exclusion needs R1. K.\\
\file{eq:etadef} & bad & Definitions agree with smx; unrestricted necessity claim does not (R6).\\
\file{eq:Sgamma} & bad & Correct baseline candidate algebra on consistent branches; candidate status (d). Executed paper 5a and smx.\\
\file{eq:BCpoles} & bad & Correct candidate poles $N\ge1$, no zero-level Gamma pole; S tests full residues. Exact spectrum remains (d).\\
\file{eq:Agoodbr} & good & Correct bound and sign boundary; subsequent root witness wrong (R5). S.\\
\file{eq:ringBA} & both & Conditional directional derivative and summed twist agree with smx, including self term; R independently differentiates classical phase. Quantum holonomy (d).\\
\file{eq:dilogres} & both & Correct audited diagonal residues, restricted constant-coefficient nine-ratio basis. No universal exclusion.\\
\file{eq:thetares} & bad & Correct on stipulated compatible analytic crossing branches and loop expansion; extend to good using its angle dictionary, not $i\pi$. R6.\\
\bottomrule
\end{longtable}
\normalsize

\C The target sections have no additional unlabelled display environments.
Important inline formulas were also checked: defect threshold/asymptotic
phase/holonomy examples; collective null vector, kinetic lower bound and
heavy-mode profiles; cutoff and sound expansions; opposite-edge identity;
the two-term convolution and its contact; Gamma Stirling/CDD algebra;
ring self derivative and composite charges; and the restricted mixed
dilogarithm argument. The failures among these are R1--R7 and R9.
The corrected $+3e_0/y^2$ term in $C_E$ is adopted correctly.
The good/bad stable-band distinction and stationary group-velocity resonance
are also transcribed correctly; the resonance location quoted is its
small-amplitude limiting value, not a formula at arbitrary $v$.

\section{Section 10 and optional improvements}
\C The order agrees with the PI's re-ranked \file{phase2/report.tex}:
good renormalisation, good two-body phase, resonant/number-changing channels,
quantum defect boundary theory, then candidate/ring selection; V1--V3 are
validation tasks. No re-ranking is requested. The good acoustic coefficient
and its scheme dependence are consistent with Section 7, and the ring
holonomy remains open consistently with Section 9.
V1's closed form is explicitly conditional and audit-pending, as required.
It must remain a \emph{directional Toda} coefficient, distinct from a
physical scalar phase and from a Gamma-relative continuum remainder.
R1, R7 and R8 repair the remaining status/definition problems without
changing the PI's scope.

\begin{itemize}
\item \C Round the numerical uncertainty upward to $8.24\times10^{-6}$,
as in the coordinator digest and errata; the displayed components sum to
$8.2298\times10^{-6}$ and a relative envelope slightly larger than
$5\times10^{-5}$. Say ``about'' that relative precision.
\item \C State $0<r<1$ immediately beside Section 10 V1's formula.
The reverse ordering follows from the separately defined exchange operation;
the $r\to1$ and $r\to0$ limits require their own control. The conditional
coefficient is not an unrestricted global analytic continuation.
\item \C Distinguish the many uses of $\eta$ (defect orientation, fluctuation,
Fredholm energy weight, Gamma argument), $h$ (Planck constant or defect
primitive), and $F$ (energy coefficient, momentum function, backflow).
R4 demonstrates an actual error caused by derivative-name migration.
\item \D Replace the long abstract's inventory with a smaller set of
qualified conclusions after the mathematical corrections. This is editorial,
not a prerequisite for retaining the independently verified coefficients.
\end{itemize}

\section{Appendix A: script-to-claim inspection ledger}
\C Every one of the 64 cited paths exists and parses as Python. The table
below records source coverage, not fresh executions of all 64 scripts.
The four paper scripts were freshly run successfully; the other entries
were inspected for their assigned calculations and compared with controlling
audits. ``Limited'' means the stated package check is narrower than the
paper claim; use the audit alongside it. A script can correctly test one
claim and fail to test another. Grouping in Appendix A must not erase that
distinction. Mathematical formulas in this review have direct script
evidence; no source inspection is represented as a new physical proof.

\small
\begin{longtable}{@{}L{6.0cm}L{9.2cm}@{}}
\toprule Script & Coverage and limitation (c)\\\midrule\endhead
\file{corrigan-zambon/defect_check.py} & Action/sewing algebra and real type-I obstruction; applies to the stated ansatz.\\
\file{corrigan-zambon/defect_transmission_phase.py} & Root transmission, delay derivatives, vacuum balances and defect symplectic term; tree level.\\
\file{corrigan-zambon/defect_two_soliton_symplectic.py} & Cross-ratio, shifts and symplectic identities; some resonance diagnostics only. Paper checks supplement these.\\
\file{corrigan-zambon/defect_vacuum_stability.py} & Sampled determinant roots and growth counts; not off-axis spectral completeness.\\
\file{corrigan-zambon/verify_defect_solutions.py} & Sampled sewing residuals for explicit transmitting solutions.\\
\file{paper/scripts/paper_conventions_check.py} & Fresh PASS. Sign dictionaries, charges, gradients and crossing samples. Historical diagnostic names are not current-source audits.\\
\file{paper/scripts/paper_draft2_checks.py} & Fresh PASS. Prior corrections, canonical signs, dictionaries and classical identities; no quantum-completeness proof.\\
\file{paper/scripts/paper_draft3_checks.py} & Fresh PASS. Many genuine algebraic tests, but important gaps in checks 1b, 1c, 2b, 4a, 4d, 5a, 5d and 5e: R10.\\
\file{paper/scripts/review_draft1_checks.py} & Fresh PASS. Independent classical review checks; optional source-suite mode was not invoked here.\\
\file{phase0/checks/convention_check.py} & Two actual symbolic assertions; no final success marker (legacy exception).\\
\file{phase1/scripts/charges.py} & Hamiltonian/GD charges, energy and velocity identities; paper checks supply the explicit vacuum-subtraction dictionary.\\
\file{phase1/scripts/figures.py} & Generator only: no scientific assertions; marker is DATA GENERATED (no checks).\\
\file{phase1/scripts/hirota_solitons.py} & Active repaired sech/PDE and finite rational/complex Hirota checks. The reported current or-True bypass is stale.\\
\file{phase1/scripts/interaction_regimes.py} & Sampled good/bad interaction and breather inequalities, resonance and shift tests; not a general quantum spectrum.\\
\file{phase1/scripts/lax_zero_curvature.py} & Symbolic Lax compatibility and explicit real/imaginary-time dictionaries.\\
\file{phase1/scripts/toda_limit.py} & Classical Toda expansion and bad-equation map, with finite-amplitude/scaling samples.\\
\file{phase2/scripts/audit_charges_action.py} & Independent charge/Legendre and regulated-action checks; endpoint and boundary-variation qualifications remain.\\
\file{phase2/scripts/audit_defq.py} & Controlling defect audit: fresh sewing action, genuine absorption, trace contractions, roots and UV checks.\\
\file{phase2/scripts/audit_det_projected.py} & Controlling geometry/trace audit: field pairings, projector tails and conditional charges. Its matrices contradict the draft.\\
\file{phase2/scripts/audit_frame_dictionary.py} & Endpoint and Jacobian checks survive. Its source/minor argument does not prove the physical exclusion: R1.\\
\file{phase2/scripts/audit_frames_defects.py} & Independent nonlinear sewing determinant and admissibility/branch-switch tests; field-frame qualifications.\\
\file{phase2/scripts/audit_onebody.py} & Known FAIL at line 383, documented rather than rerun. Several interpretations are superseded by the scheme audit.\\
\file{phase2/scripts/audit_onebody_scheme.py} & Controlling vacuum/one-body audit; independent regularisations, refined edge checks, and diagnosis of the earlier failed tolerance.\\
\file{phase2/scripts/audit_phase2i.py} & Controlling symbol, transform/reversion, convolution and finite-data audit; uniform limits and physical frame explicitly open.\\
\file{phase2/scripts/audit_projected_measure.py} & Actual collective field geometry and collision quadrature; projector endpoint failure, not a full determinant.\\
\file{phase2/scripts/audit_smx.py} & Controlling candidate audit: pole residues, ring derivatives and direct crossing quadrature, with restricted exclusions.\\
\file{phase2/scripts/audit_toda_dictionary.py} & Fresh finite-sea endpoint/phase calculations and classical-frame obstruction; physical quantum map open.\\
\file{phase2/scripts/charge_and_frame_audit.py} & Charge-pair algebra and numerical frame diagnostics; no physical determinant identification.\\
\file{phase2/scripts/consistency_checks.py} & Multiple classical identities; bad crossing and some asymptotic comparisons are printed only and need paper-check supplementation.\\
\file{phase2/scripts/defect_smatrix_good_continuation.py} & Candidate/fusion tests, including oriented zero-constant tree fusion; not quantum boundary normalisation.\\
\file{phase2/scripts/defq_absorption_fusion.py} & Sampled holonomies and momentum remainders, not a universal no-go; controlling defq audit qualifies interpretation.\\
\file{phase2/scripts/defq_classical_charges.py} & Classical charge algebra and transmission tests; original negative-orientation absorption input is wrong. Use audit\_defq.\\
\file{phase2/scripts/defq_oneloop_renorm.py} & Determinant, UV and operator algebra; energy-sign/asymptotic printed descriptions require audit\_defq.\\
\file{phase2/scripts/det_charges.py} & Trace/series and kernel calculations; some Legendre/frame checks are tautological. Use audit\_det\_projected for independent evidence.\\
\file{phase2/scripts/det_geometry.py} & Instantaneous geometry and moments; abstract projector test does not test physical endpoint tails.\\
\file{phase2/scripts/det_spectrum.py} & Mode calculations; formal pencil check includes tautologies. Audit supplies differentiated modes and completeness qualifications.\\
\file{phase2/scripts/frame_analysis.py} & Saved-data fits, error sensitivity and mixed-remainder diagnostics; no physical phase selection.\\
\file{phase2/scripts/frame_edge.py} & Finite-h solver scan at inverted momenta; saved numerical evidence, no uniform-limit theorem.\\
\file{phase2/scripts/frame_identities.py} & Finite-sea endpoint tests and algebraic counterexample; generic fixed-momentum chain rule is tautological, repaired by audit.\\
\file{phase2/scripts/onebody_continuum.py} & Good/formally continued/stable-band traces and asymptotic distinctions; physical bad determinant not established.\\
\file{phase2/scripts/onebody_prec_analysis.py} & Unconstrained fits and deterministic sensitivity envelopes from saved data; central value is not a rigorous error bound.\\
\file{phase2/scripts/onebody_prec_edge.py} & Corrected-transform inversion and conditional edge series/energy coefficients.\\
\file{phase2/scripts/onebody_prec_identities.py} & Finite-sea dispersion derivative identity; default run validates cached residuals unless recompute is requested.\\
\file{phase2/scripts/onebody_prec_scan.py} & Refined finite-h momentum/energy scan and mesh checks; writes data, not rerun here.\\
\file{phase2/scripts/onebody_scheme_continuum.py} & Vacuum strain energies, finite quadratic matching, soliton moments and conditional trace slopes.\\
\file{phase2/scripts/onebody_scheme_dictionary.py} & Unit factors, thermodynamic identities and harmonic Toda vacuum shift.\\
\file{phase2/scripts/routes_BC_checks.py} & Affine Toda factorisation has numerical assertions; bracket algebra is supplemented by review smeared-bracket tests.\\
\file{phase2/scripts/semiclassical_phase.py} & Classical phase-gradient identity; quantum candidate necessity does not follow.\\
\file{phase2/scripts/smatrix_ansatz.py} & Candidate/root dictionary and phase algebra, with numerical conditions; ansatz remains conjectural.\\
\file{phase2/scripts/smatrix_poles_bootstrap.py} & Full candidate breather residue and fusion-product tests; not an exact physical-spectrum proof.\\
\file{phase2/scripts/smatrix_rejected_forms.py} & Finite root-difference least-squares basis and mirror mixed derivatives; no unrestricted Gamma-prefactor exclusion.\\
\file{phase2/scripts/smatrix_semiclassical_unitarity_crossing.py} & Numerical candidate unitarity, continuation and Stirling checks; cannot prove physical quantum crossing.\\
\file{phase2/scripts/smx_dictionary_and_families.py} & Family algebra and fits; broad exclusion proof is refuted. Use audit\_smx restricted residue proof.\\
\file{phase2/scripts/smx_poles_good_vs_bad.py} & Candidate pole samples; original sign/positivity claims are not tested. Use audit\_smx and R5.\\
\file{phase2/scripts/smx_residual_ambiguity_interface.py} & Original crossing offset is inserted into J, not obtained by independent quadrature; use audit\_smx.\\
\file{phase2/scripts/smx_twisted_ring_bethe.py} & Assumes seam twist and substitutes pair derivative; use audit\_smx for independent gradients and self derivative.\\
\file{phase2/scripts/symplectic_scattering.py} & Bad-equation symplectic identities with and without cross terms.\\
\file{phase2/scripts/toda_ba_vs_label.py} & Sampled Bethe-versus-label phase derivative ratios; classical dictionary only.\\
\file{phase2/scripts/toda_bethe_classical.py} & Classical dressed momentum/energy versus soliton charges.\\
\file{phase2/scripts/toda_dressed_phase_check.py} & Dressed phase derivatives versus classical shifts; not quantum asymptotic completeness.\\
\file{phase2/scripts/toda_label_symplectic.py} & Lattice soliton canonical pairings and label-frame symplectic identities.\\
\file{phase2/scripts/toda_scaling_fast.py} & Printed Boolean-gated success for continuum scaling; no assertion/nonzero exit on failure.\\
\file{phase2/scripts/toda_soliton_classical.py} & Lattice soliton mass, momentum and energy via sampled lattice sums.\\
\file{phase2/scripts/toda_wh_and_kp_synthesis.py} & Conditional coefficient integrals/reversion. Printed symbol/contact/reference and physical claims exceed actual tests; use phase2i audit.\\
\bottomrule
\end{longtable}
\normalsize

\section{Runs and failures recorded}
\C Fresh runs from \file{paper/scripts/}, with bytecode writing disabled,
completed with exit code 0 and printed \file{ALL CHECKS PASSED}:
\file{paper_conventions_check.py}, \file{paper_draft2_checks.py},
\file{paper_draft3_checks.py}, and \file{review_draft1_checks.py}.
The independent \file{review_draft3_checks.py} also prints
\file{ALL CHECKS PASSED}. Its success means the corrected identities and
the identified counterexamples passed, not that the draft is error-free.
The supplied PDF is built twice with \file{pdflatex}; auxiliary files and
rendered QA images are kept outside the repository.

\begin{enumerate}[leftmargin=*]
\item \C Draft failures reproduced: opposite matrix signs; erroneous
projector source definitions; loop-order inconsistency and reversed finite
tadpole sign; wrong four-term source and unweighted coordinate test; wrong
contact derivative; and false tau-zero witness. None is hidden by a passing
determinant, minor, or symbolic cancellation check.
\item \C Unsupported scope claims are retained as failures of certification:
physical rank-two reduction, exact crossing as a necessity, unrestricted
Gamma-argument exclusion, stability completeness, universal formula coverage,
and stale current-status statements. The corresponding narrower algebraic
or conditional results are distinguished in the correction list.
\item \C The earlier \file{audit_onebody.py} line-383 failure and Phase-2d
successive-step-bound failure are preserved as recorded in the errata and
inventory. They were not rerun here because doing so can write existing
files. The true limiting one-body identity is not discredited by the crude
extrapolation's failed tolerance.
\item \C Review-check development: the first run compared the exact
inverse-transform Krein source against the later package's rounded
\emph{finite-difference convolution derivative} with tolerance $10^{-17}$.
It failed by $2.0655\times10^{-17}$, matching the package's explicitly
reported $2.07\times10^{-17}$ differentiation error. That comparison now
uses $2.1\times10^{-17}$ and is named as a finite-difference-error comparison;
an independent direct differentiation of the inverse transform is additionally
required at $10^{-48}$. This tolerance change is disclosed, not presented
as an exact equality or a repaired source formula.
\item \C A second development run reached the final source-status check
and failed because the review searched for a bare numeric threshold string,
whereas the original assertion wraps the threshold in \file{mp.mpf}.
The search was corrected to the actual assertion text; no numerical
criterion or original file was changed.
\item \C The first PDF build exposed an overfull line containing the
unbroken SHA-256 digest. It was split across two lines before the final
two-pass build. During prose QA, an intermediate sentence incorrectly said
both proposed projector inputs vanish on $\eta=(w,0)$; only the proposed
$b_r$ does. The sentence was corrected; the independent constant-tail
counterexample and all mathematical checks are unchanged.
One editorial patch attempt used mismatched context and applied no change;
the intended addition was then applied against the actual text.
\end{enumerate}

\paragraph{Reproduction.}
\C Run \file{python3 -B paper/scripts/review_draft3_checks.py}.
Fresh Fredholm checks use 50 digits and 40 Gauss nodes with tolerance
$10^{-27}$; their finite quadrature accuracy is not claimed to equal the
working precision. Gamma residue checks compare displacements
$10^{-15}$ and $10^{-20}$ with tolerance $10^{-12}$.
WH coefficient checks use the explicitly conditional bound
$|\mathrm{residual}|y^3<0.0004$ at two independent inversion points;
they supplement the source's exact formal reversion, not a uniform theorem.
All these tolerances are fixed in the checker.
\D The physical two-body phase, uniform limits, full determinant,
boundary state space and exact quantum spectrum remain unresolved.
\end{document}
